\section{Resolved difficulties and the remaining mathematical boundary}\label{sec:resolved}

The framework provides more than a uniform instruction to supply moduli.
Several of the required data are \emph{derived} from existing presentations:
a dominating $L^1$ name supplies a tail cutoff; a regular simple name
supplies a conditional-expectation convergence rate; a compact net and
distance-function names supply continuous tent quadrature; a contraction
and a perturbation gap supply inverse and homotopy clocks. Each derivation
has an explicit finite routine and a proof of its numerical contract.

\begin{longtable}{@{}p{0.28\textwidth}p{0.40\textwidth}p{0.25\textwidth}@{}}
\toprule
Difficulty & Resolution obtained & Remaining hypothesis\\
\midrule
\endhead
Incorrect separate-tail truncation in the earlier Vitali proof & Replace it by the checked clipping estimate, then build the limit from pairwise Cauchy data. & A uniform cutoff program and a capped-distance Cauchy rate.\\
The limit is assumed integrable before convergence is proved & The diagonal simple-code sampler constructs the $L^1$ limit and derives its uniform tails. & Joint PR names for the sequence.\\
An additional amplitude-tail oracle appears necessary under domination & Extract a cutoff from the finite supremum of each simple approximation to the dominator. & A regular $L^1$ name for the dominator.\\
Small cells cause large posterior values and unstable division & Use the full $L^1$ conditional operator; the cell mass cancels the reciprocal. & A positive mass margin for a pointwise posterior query.\\
Finite operators fail to approximate the identity uniformly & Compose strong input-dependent actions with certified gain schedules. & An action routine on every supplied dense input; no automatic operator-name extraction.\\
Integration is absent from compact model evaluation & Add rational tent quadrature and an expectation binder with weighted errors. & Compact nets, compatible measure data, Lipschitz/range bounds.\\
An analytic integral has a singular spatial factor & Use $L^1$ coefficient bounds and integrate the norm-convergent series. & A radius gap and a jointly presented coefficient family.\\
Ordinary nerve methods require acyclic overlaps & Retain local chain models and transfer their full differential by finite homological perturbation. & Explicit local models and a valid cover-to-total-complex comparison.\\
An approximate differential is mistakenly treated as square-zero & Carry the exact residual identities and a geometric bound until the completed differential or exact nilpotent formula is available. & Exact square-zero of the original perturbed differential.\\
\bottomrule
\end{longtable}

These hypotheses cannot all be erased. The bounded martingale construction
in Section~\ref{sec:measure} makes a uniform computable convergence-rate
extractor decide halting. Small $L^1$ error cannot control a supremum.
Independent fair events prevent a compact-net treatment of an entire
atomless probability algebra. Small differential perturbations can change
homology; the Schur-complement example shows the change on an explicit
retained model. Exact real equality, unbounded disintegration, arbitrary
model selection, and unrestricted countable unions remain outside the
proved construction unless their missing data are supplied or derived.

For further repository work, the highest-value next object-language step
is the integer clipping inequality \eqref{eq:object-vitali}, followed by
the regular-name diagonal and conditional-average contraction. They use
the existing arithmetic backend and produce reusable continuous mathematics
with a clear verification boundary. The matrix residual endpoint is a
parallel target with similarly explicit bounded-sum invariants. Finishing
these endpoints would turn the present full mathematical arguments and
checked finite kernels into new, inspectable \texttt{Prf PRAAx} exports.
